% W4i43_Observables.tex
% DFD 17-Agent Research Programme --- Iteration 43 (i43)
% Agents 8 (Laboratory), 9 (Particle Physics), 10 (Astrophysics), 11 (CMB), 12 (LSS), 13 (Dark Sector), 14 (SymBoltz/Numerics)
% Theme: Experimental updates, Th-229 progress, DESI, SymBoltz publication, epsilon_H clarification
% Date: 2026-03-23

\documentclass[11pt,a4paper]{article}
\usepackage[utf8]{inputenc}
\usepackage{amsmath,amssymb,amsfonts,amsthm}
\usepackage{hyperref}
\usepackage{booktabs}
\usepackage{longtable}
\usepackage{geometry}
\usepackage{xcolor}
\usepackage{tcolorbox}
\geometry{margin=2.5cm}

\newtheorem{theorem}{Theorem}
\newtheorem{proposition}{Proposition}
\newtheorem{lemma}{Lemma}
\newtheorem{corollary}{Corollary}
\newtheorem{definition}{Definition}
\newtheorem{remark}{Remark}

\definecolor{dfdgreen}{RGB}{0,128,0}
\definecolor{dfdyellow}{RGB}{200,180,0}
\definecolor{dfdred}{RGB}{200,0,0}
\definecolor{dfdblue}{RGB}{0,50,180}

\newcommand{\GREEN}{\textcolor{dfdgreen}{\textbf{GREEN}}}
\newcommand{\YELLOW}{\textcolor{dfdyellow}{\textbf{YELLOW}}}
\newcommand{\RED}{\textcolor{dfdred}{\textbf{RED}}}
\newcommand{\NEW}{\textcolor{dfdblue}{\textbf{NEW}}}
\newcommand{\RESOLVED}{\textcolor{dfdgreen}{\textbf{RESOLVED}}}
\newcommand{\Tr}{\operatorname{Tr}}
\newcommand{\CP}{\mathbb{C}P}

\title{DFD i43: Observables --- Laboratory, Particle, Astro, CMB, LSS, Dark Sector, Numerics\\[6pt]
\large Agents 8, 9, 10, 11, 12, 13, 14\\[6pt]
\normalsize Iteration 12/20 of Quantum Gravity Cycle (i32--i51)}
\author{DFD 17-Agent Research Programme}
\date{2026-03-23}

\begin{document}
\maketitle
\tableofcontents
\newpage

%=============================================================================
\part{Agent 8: Laboratory Experiments}
\label{part:agent8}
%=============================================================================

\section{Mandate}

Update on laboratory experiments relevant to DFD: Th-229 nuclear clock, SRF cavities, SQMS, and precision measurements.

\section{\NEW: Thorium-229 Nuclear Clock Breakthroughs (2025--2026)}

Major progress has occurred in the Th-229 nuclear clock programme, which is critical for DFD's $\dot{\alpha}/\alpha$ prediction at $10^{-19}$/yr:

\begin{enumerate}
\item \textbf{Frequency reproducibility} (Nature, January 2026): The $^{229}$Th:CaF$_2$ nuclear clock transition has been characterized for frequency reproducibility as a function of doping concentration, temperature, and time. This is a key step toward operational nuclear clocks.

\item \textbf{Temperature sensitivity} (PRL 134, March 2025): Temperature-induced frequency shifts characterized. To reach $10^{-18}$ fractional precision, crystal temperature must be stabilized to within $5~\mu$K. This sets the engineering challenge for DFD-relevant sensitivity.

\item \textbf{Fine-structure constant sensitivity} (Nature Communications, 2025): Ma Yugang's team at Fudan University computed the nuclear isotope shifts and predicted the absolute transition frequency of isolated $^{229}$Th$^{3+}$ using non-perturbative multiconfiguration Dirac--Hartree--Fock. The $\alpha$-sensitivity coefficient $K_\alpha$ has been determined.

\item \textbf{Internal conversion electron M\"ossbauer spectroscopy} (December 2025): Eric Hudson's UCLA team achieved laser-driven measurement of $^{229}$ThO$_2$. A new detection method that could accelerate development.

\item \textbf{Frequency ratio measurement} (Nature, 2024): The frequency ratio of $^{229\text{m}}$Th to the $^{87}$Sr atomic clock has been measured, establishing the nuclear transition in the metrological framework.
\end{enumerate}

\subsection{Timeline revision}

Previous estimate: Th-229 sensitivity to $\dot{\alpha}/\alpha$ at $10^{-19}$/yr by $\sim$2030. With the recent frequency reproducibility demonstration:
\begin{itemize}
\item Operational solid-state Th-229 clock: $\sim$2028
\item $10^{-18}$ fractional sensitivity: $\sim$2029
\item DFD-relevant $10^{-19}$/yr sensitivity: $\sim$2030--2031 (unchanged)
\end{itemize}

\section{SQMS and SRF Cavities}

No new results since i42. SQMS 2.0 renewal (\$125M) continues. Dark SRF Phase 1 advancing.

\section{Agent 8 Assessment: Grade A-}

Significant Th-229 developments strengthen the DFD $\dot{\alpha}/\alpha$ test timeline. Multiple independent groups making progress.

\section{New Literature (Agent 8)}

\begin{enumerate}
\item \textbf{Nature (January 2026)} --- Frequency reproducibility of solid-state $^{229}$Th nuclear clocks.
\item \textbf{PRL 134 (March 2025)} --- Temperature sensitivity of Th-229 clock transition.
\item \textbf{Nature Comms (2025)} --- Fine-structure constant sensitivity of Th-229 transition.
\item \textbf{Hudson et al.\ (December 2025)} --- IC electron M\"ossbauer spectroscopy of $^{229}$ThO$_2$.
\item \textbf{Nature (2024)} --- $^{229\text{m}}$Th/$^{87}$Sr frequency ratio measurement.
\item \textbf{Nuclear Physics News (2025)} --- Review: toward a Th-229 nuclear clock.
\end{enumerate}


%=============================================================================
\part{Agent 9: Particle Physics}
\label{part:agent9}
%=============================================================================

\section{Mandate}

Update on muon $g-2$ final result, LHC Run 3, and particle physics tests of DFD.

\section{Muon $g-2$: Final Result Confirms GREEN Status}

The Fermilab Muon $g-2$ experiment released its final result on June 3, 2025:

\begin{equation}
a_\mu^{\text{exp}} = 0.001\,165\,920\,705(114)
\end{equation}

Key points:
\begin{itemize}
\item The experimental value deviates by $5.1\sigma$ from the 2020 WP theory prediction (dispersive HVP).
\item However, it differs by only $\sim 1\sigma$ from the lattice QCD prediction (BMW collaboration and subsequent confirmations).
\item The lattice-based value has been adopted as the representative SM prediction.
\item \textbf{DFD prediction}: $g-2$ consistent with SM (no BSM $\psi$-field contribution). This is \GREEN.
\end{itemize}

\textbf{Status}: Promoted to \GREEN{} in i42. Confirmed. The tension is between theory methods (dispersive vs.\ lattice), not between experiment and theory.

\section{LHC Run 3: No BSM Signals}

LHC Run 3 continues with no BSM signals. This is consistent with DFD, which predicts no new particles beyond the SM (all new physics is geometric/topological, not in the particle spectrum).

\section{Agent 9 Assessment: Grade B}

$g-2$ status confirmed. No new particle physics developments affect DFD.

\section{New Literature (Agent 9)}

\begin{enumerate}
\item \textbf{Fermilab $g-2$ final} (June 2025) --- $a_\mu = 0.001165920705(114)$; lattice QCD agreement.
\item \textbf{BMW lattice update (2025)} --- Confirmed $a_\mu^{\text{SM,lattice}}$ consistent with experiment.
\item \textbf{LHCb Run 3 (2025--2026)} --- Lepton universality tests; anomalies dissolved.
\item \textbf{ATLAS/CMS Higgs (2025)} --- $m_H = 125.11 \pm 0.11$ GeV; DFD compatible.
\item \textbf{CERN Courier (2025)} --- Overview of $g-2$ resolution.
\end{enumerate}


%=============================================================================
\part{Agent 10: Astrophysics}
\label{part:agent10}
%=============================================================================

\section{Mandate}

Wide binary update. Astrophysical tests of gravity modifications.

\section{Wide Binaries: Field Remains Divided}

The wide binary test of gravity in the low-acceleration regime ($a < a_0 \approx 1.2 \times 10^{-10}$ m/s$^2$) continues to produce conflicting results:

\begin{itemize}
\item \textbf{Cookson et al.\ (2026)} (arXiv:2504.07569): Using Gaia DR3 with improved triple-star modelling, find Newtonian gravity significantly preferred over MOND. Updated selection cuts reduce contamination.
\item \textbf{Hernandez et al.\ (2024)}: Critical review arguing that implementations finding Newtonian preference share statistical design problems that bias conclusions.
\item \textbf{Banik et al.\ (2024)} (MNRAS 533): Report MOND-consistent signal in Gaia data; contest the Newtonian preference.
\item \NEW{} \textbf{arXiv:2603.11015} (March 2026): ``Gravitational Anomaly Measurement in Wide Binaries is Sensitive to Orbital Modeling.'' Shows that the result depends sensitively on how orbits are modelled, explaining some of the disagreement.
\item \NEW{} \textbf{arXiv:2512.25002} (December 2025): Distributions of wide binary stars; orbital momenta, masses, and MOND manifestations. Finds possible MOND signature in orbital momentum distributions.
\end{itemize}

\textbf{DFD position}: DFD predicts no MOND-like departure from Newtonian gravity in wide binaries (the $\psi$-field EFE is at a different scale). The Cookson et al.\ result is consistent with DFD, but the field is not settled. Status remains at $3\sigma$, \YELLOW.

\section{Agent 10 Assessment: Grade B}

Wide binary field remains active and divided. New orbital modelling paper may help resolve discrepancies.

\section{New Literature (Agent 10)}

\begin{enumerate}
\item \textbf{Cookson et al.\ (2026)}, arXiv:2504.07569 --- Gaia DR3 wide binaries; GR preferred.
\item \textbf{arXiv:2603.11015} (March 2026) --- Orbital modelling sensitivity in wide binaries. \NEW.
\item \textbf{arXiv:2512.25002} (December 2025) --- Wide binary orbital momenta and MOND. \NEW.
\item \textbf{Hernandez et al.\ (2024)} --- Critical review of Gaia wide binary gravity tests.
\item \textbf{Banik et al.\ (2024)}, MNRAS 533 --- Counter-argument for MOND signal.
\end{enumerate}


%=============================================================================
\part{Agent 11: CMB}
\label{part:agent11}
%=============================================================================

\section{Mandate}

$A_L$ update. Simons Observatory status. $\varepsilon_H$ clarification (the $n_s$ concern from Agent 16).

\section{$A_L$ Lensing: Demoted to YELLOW (from i42)}

The joint ACT+SPT+Planck measurement:
\begin{equation}
A_L^{\text{recon}} = 1.025 \pm 0.017
\end{equation}
from Qu et al.\ (PRL 136, January 2026; arXiv:2504.20038). The SNR of the combined lensing bandpowers is 61.

DFD's previous analytic estimate of $A_L \sim 1.10$--$1.20$ is in $\sim 4\sigma$ tension. Revised to $A_L^{\text{DFD}} \sim 1.00$--$1.05$ pending Boltzmann computation. Status: \YELLOW.

\section{\RESOLVED: $\varepsilon_H = 0.05$ Is NOT the Slow-Roll Parameter}
\label{sec:epsilon_H_resolution}

\begin{tcolorbox}[colback=green!5!white, colframe=green!60!black,
  title={RESOLUTION: $\varepsilon_H$ vs.\ $\epsilon_V$ Confusion}]

Agent 16 in i42 flagged that $\varepsilon_H = 0.05$ gives $n_s = 1 - 2\varepsilon_H = 0.90$, which is $>10\sigma$ from the observed $n_s = 0.965 \pm 0.004$.

\textbf{This is a category error.} The quantity $\varepsilon_H = N_{\text{gen}} / k_{\max} = 3/60 = 0.05$ is the \textbf{Higgs localization width}---the squared overlap $|\langle i | H \rangle|^2$ in the channel Hilbert space. It is \textbf{not} the inflationary slow-roll parameter $\epsilon_V \equiv (M_P^2/2)(V'/V)^2$.

The relation $n_s = 1 - 2\epsilon_V - \eta_V$ applies to the slow-roll parameter $\epsilon_V$, not to $\varepsilon_H$. These are different physical quantities:

\begin{center}
\begin{tabular}{lll}
\toprule
\textbf{Symbol} & \textbf{Definition} & \textbf{Physical meaning} \\
\midrule
$\varepsilon_H$ & $N_{\text{gen}} / k_{\max} = 3/60$ & Higgs localization width \\
$\epsilon_V$ & $(M_P^2/2)(V'/V)^2$ & Slow-roll parameter \\
\bottomrule
\end{tabular}
\end{center}

The DFD inflationary sector would need its own slow-roll analysis, which is a separate computation not yet performed. The $\varepsilon_H = 0.05$ value has \textbf{no direct implication} for $n_s$.
\end{tcolorbox}

\section{\NEW: Simons Observatory Update}

\begin{itemize}
\item The Large Aperture Telescope (LAT) achieved first light in February 2025, capturing an image of Mars.
\item Three Small Aperture Telescopes (SATs) have been observing since mid-2024.
\item \NEW: Additional telescopes funded by Japan and the UK are scheduled for 2026, doubling the number of SATs.
\item First CMB science results expected: 2026--2027.
\item DFD-relevant: SO will provide independent confirmation of $A_L^{\text{recon}}$ and improve $n_s$ precision.
\end{itemize}

\section{Agent 11 Assessment: Grade A-}

Two critical resolutions: $\varepsilon_H$ category error clarified; $A_L$ status updated. SO making good progress.

\section{New Literature (Agent 11)}

\begin{enumerate}
\item \textbf{Qu et al.\ (2026)}, PRL 136; arXiv:2504.20038 --- Joint lensing; $A_L = 1.025 \pm 0.017$.
\item \textbf{Simons Observatory LAT first light} (February 2025) --- Mars observation.
\item \textbf{SO expansion} (2026) --- Japan/UK-funded additional SATs.
\item \textbf{ScienceDaily (March 2026)} --- Cosmic birefringence signal update from CMB polarization.
\item \textbf{CMB-S4 forecasts} (arXiv:2505.22775, May 2025) --- Constraints on non-thermal light relics.
\end{enumerate}


%=============================================================================
\part{Agent 12: Large-Scale Structure}
\label{part:agent12}
%=============================================================================

\section{Mandate}

DESI DR2 update. Euclid DR1 timeline. LSS predictions.

\section{DESI DR2: Dynamical Dark Energy Debate Continues}

\begin{itemize}
\item DESI DR2 BAO measurements cover $0.1 < z < 4.2$ using the Lyman-$\alpha$ forest and galaxies.
\item Combined with CMB and SN data: preference for dynamical DE ($w_0w_a$CDM over $\Lambda$CDM) at $2.8$--$4.2\sigma$ depending on SN dataset.
\item \NEW: Giare et al.\ (arXiv:2504.15222, April 2025; published EPJC): ``Did DESI DR2 truly reveal dynamical dark energy?'' --- Argues that tensions among CMB, BAO, and SN datasets make the claim not robust.
\item \NEW: Robust evidence paper (Science Direct, 2026): Joint DESI DR2 + ACT + SPT + Planck analysis finds robust evidence for dynamical DE, especially at low redshift ($z < 0.3$).
\item The $N_{\text{eff}} = 3.23^{+0.35}_{-0.34}$ constraint is consistent with SM (and DFD) expectations.
\end{itemize}

\textbf{DFD position}: DFD predicts an effective dynamical DE signature from the distance bias, not from a physical DE equation of state. The DESI signal, if confirmed, would be consistent with DFD's distance-bias mechanism. Status: \YELLOW{} (the signal is debated).

\section{Euclid DR1 Timeline}

\begin{itemize}
\item \textbf{Euclid DR1}: Expected October 2026. Covers $\sim$2300 deg$^2$ (1 year of survey data).
\item Quick Data Release 1 (Q1) already released with early science results and images.
\item 26 million galaxies imaged; most distant at $z > 2$.
\item DFD pre-registration deadline: 20 June 2026 (on track).
\item 7 DFD observables identified for pre-registration (from i41).
\end{itemize}

\section{Agent 12 Assessment: Grade B+}

DESI dynamical DE debate sharpens. Euclid DR1 timeline confirmed for October 2026.

\section{New Literature (Agent 12)}

\begin{enumerate}
\item \textbf{Giare et al.\ (2025)}, arXiv:2504.15222; EPJC --- Challenges DESI DR2 dynamical DE claim. \NEW.
\item \textbf{Robust evidence paper (2026)} --- Joint DESI+ACT+SPT+Planck supports dynamical DE.
\item \textbf{DESI DR2 Guide} (March 2025) --- Official results summary.
\item \textbf{Euclid Q1 release} (2025) --- Quick data release with early images.
\item \textbf{JWST + DESI DR2} (Phys.\ Rev.\ D, 2025) --- Measuring neutrino masses jointly.
\end{enumerate}


%=============================================================================
\part{Agent 13: Dark Sector}
\label{part:agent13}
%=============================================================================

\section{Mandate}

Neutrino mass constraint update. Dark matter status.

\section{$\Sigma m_\nu$: Under Pressure (\YELLOW)}

\begin{center}
\begin{tabular}{lcc}
\toprule
\textbf{Constraint} & \textbf{Upper bound} & \textbf{DFD prediction} \\
\midrule
DESI DR2 + CMB (Bayesian 95\%) & $< 0.064$ eV & $\approx 0.061$ eV \\
DESI DR2 + CMB (frequentist 95\%) & $< 0.053$ eV & $\approx 0.061$ eV \\
$w_0 w_a$CDM relaxation & $< 0.163$ eV & $\approx 0.061$ eV \\
\bottomrule
\end{tabular}
\end{center}

\textbf{Assessment}: DFD's prediction of $\Sigma m_\nu \approx 0.061$ eV is:
\begin{itemize}
\item Safe under the Bayesian $\Lambda$CDM bound ($0.061 < 0.064$)
\item Excluded at $1.1\sigma$ under the frequentist bound ($0.061 > 0.053$)
\item Comfortable under $w_0 w_a$CDM relaxation
\end{itemize}

\NEW: arXiv:2603.10787 (March 2026) combines ACT DR6 + DESI DR2 for neutrino mass measurement. The constraint continues to tighten.

\textbf{Honest assessment}: The frequentist bound is uncomfortably close. If future data tighten the Bayesian bound below 0.055 eV, DFD would face a genuine tension. Status: \YELLOW.

\section{Dark Matter: No Change}

DFD predicts dark matter is an effective phenomenon arising from the $\psi$-field's gravitational distance bias. No direct detection is expected (consistent with all null results from XENON, LZ, PandaX).

\section{Agent 13 Assessment: Grade B}

Neutrino mass constraint updated. DFD marginally safe but honest about the pressure.

\section{New Literature (Agent 13)}

\begin{enumerate}
\item \textbf{DESI DR2 neutrino} (arXiv:2503.14744) --- $\Sigma m_\nu < 0.064$ eV.
\item \textbf{arXiv:2603.10787} (March 2026) --- ACT DR6 + DESI DR2 neutrino mass. \NEW.
\item \textbf{PRL (2025)} --- Positive neutrino masses via matter-to-DE conversion.
\item \textbf{JWST + DESI joint} (Phys.\ Rev.\ D, 2025) --- Joint neutrino mass measurement.
\item \textbf{DESI reanalysis} (arXiv:2511.20757, November 2025) --- Constraints on DE, curvature, $\nu$ masses.
\end{enumerate}


%=============================================================================
\part{Agent 14: SymBoltz and Numerical Methods}
\label{part:agent14}
%=============================================================================

\section{Mandate}

SymBoltz Phase 0 progress. New Boltzmann solver developments.

\section{\NEW: SymBoltz.jl Published in A\&A}

\begin{tcolorbox}[colback=green!5!white, colframe=green!60!black,
  title={MAJOR DEVELOPMENT: SymBoltz.jl Published}]
\textbf{SymBoltz.jl: A symbolic-numeric, approximation-free, and differentiable linear Einstein--Boltzmann solver}

Published in \emph{Astronomy \& Astrophysics} (2026, Volume A\&A, March issue).

Key features:
\begin{itemize}
\item Julia package for solving the linear Einstein--Boltzmann equations
\item Symbolic-numeric interface for specifying equations
\item Free of approximation switching schemes (unlike CLASS/CAMB)
\item Compatible with automatic differentiation
\item Ideal for non-standard cosmological models
\end{itemize}
\end{tcolorbox}

\subsection{Implications for DFD Phase 0}

SymBoltz.jl is the \textbf{exact tool} needed for DFD's SymBoltz Phase 0:
\begin{enumerate}
\item \textbf{Approximation-free}: DFD's modifications to the Boltzmann hierarchy require consistent treatment without the usual approximation switches (e.g., tight-coupling, RSA). SymBoltz.jl provides this natively.
\item \textbf{Symbolic equations}: DFD's distance-bias modification to the luminosity distance can be implemented symbolically, making the modification transparent and auditable.
\item \textbf{Differentiable}: Gradient-based optimization for DFD parameters is available out of the box.
\item \textbf{Julia ecosystem}: Compatible with existing DFD numerical tools.
\end{enumerate}

\subsection{Phase 0 Plan Update}

With SymBoltz.jl now published, the Phase 0 plan is updated:

\begin{center}
\begin{tabular}{clcc}
\toprule
\textbf{Day} & \textbf{Task} & \textbf{Hours} & \textbf{Status} \\
\midrule
1 & Install SymBoltz.jl, reproduce $\Lambda$CDM $C_\ell$ & 2 & Ready \\
1 & Implement DFD distance-bias modification & 3 & Ready \\
2 & Run $\Lambda$CDM + distance bias, compare $C_\ell^{TT}$ & 2 & Ready \\
2 & Extract $R_{31}$ third peak ratio & 1 & Ready \\
2 & Compute $A_L^{\text{eff}}$ from modified $C_\ell$ & 2 & Ready \\
3 & Go/no-go assessment & 2 & Ready \\
\bottomrule
\end{tabular}
\end{center}

\textbf{Estimated total}: 12 hours. \textbf{Go/no-go criteria} (from i41):
\begin{enumerate}
\item $R_{31}$ within $2\sigma$ of Planck? If yes, \GREEN{} (and $R_{31}$ promoted from \YELLOW).
\item $A_L^{\text{eff}}$ within $2\sigma$ of $1.025 \pm 0.017$? If yes, $A_L$ can be re-promoted.
\item $n_s$ consistent with $0.965 \pm 0.004$? If yes, spectral index OK.
\end{enumerate}

\section{\NEW: gCAMB (GPU-Accelerated CAMB)}

A GPU-accelerated version of CAMB (gCAMB) was published in 2026. Key features:
\begin{itemize}
\item Preserves all CPU-only CAMB features
\item Significantly accelerates power spectrum generation
\item Halves power consumption in high-accuracy settings
\end{itemize}

This is relevant for DFD parameter space exploration if SymBoltz.jl proves too slow for MCMC runs.

\section{Agent 14 Assessment: Grade A}

The publication of SymBoltz.jl is a breakthrough for DFD Phase 0. The tool is perfectly suited. Phase 0 execution should commence immediately.

\section{New Literature (Agent 14)}

\begin{enumerate}
\item \textbf{SymBoltz.jl} (A\&A, March 2026) --- Approximation-free, differentiable Boltzmann solver. \NEW. \textbf{Critical for DFD.}
\item \textbf{gCAMB} (arXiv:2509.25110; A\&C, 2026) --- GPU-accelerated CAMB. Backup option.
\item \textbf{CAMB 1.6.5} (February 2026) --- Latest release.
\item \textbf{Lee et al.\ (2025)} --- CLASS speedup via integral equations.
\item \textbf{CMBAns} (arXiv:1910.00725) --- Alternative $C_\ell$ solver; pedagogical.
\end{enumerate}


%=============================================================================
\section*{Summary: Observable Agents i43}
%=============================================================================

\begin{center}
\begin{tabular}{llp{9cm}}
\toprule
\textbf{Agent} & \textbf{Grade} & \textbf{Key Result} \\
\midrule
8 (Lab) & A- & Th-229 frequency reproducibility demonstrated (Nature, 2026). $\dot{\alpha}/\alpha$ test on track for 2030. \\
9 (Particle) & B & $g-2$ final confirmed \GREEN. No BSM at LHC. \\
10 (Astro) & B & Wide binaries still divided; new orbital modelling paper. \\
11 (CMB) & A- & \RESOLVED: $\varepsilon_H$ is NOT slow-roll. $A_L$ \YELLOW. SO expanding. \\
12 (LSS) & B+ & DESI dynamical DE debated. Euclid DR1 October 2026. \\
13 (Dark) & B & $\Sigma m_\nu$ under pressure; DFD marginally safe. \\
14 (Numerics) & A & SymBoltz.jl published in A\&A. Phase 0 ready to execute. \\
\bottomrule
\end{tabular}
\end{center}

\end{document}
